\section{Database Schema (PostgreSQL)}
\label{sec:implementation:database-schema}

The PostgreSQL schema is structured to represent railway infrastructure in a way that supports real-time interlocking validation, efficient querying, and strict audit compliance. It is organized into three logical namespaces: one dedicated to core control entities such as signals, track segments, and point machines; another for configuration data like signal types, aspects, and point positions; and a third reserved for audit logging of operational events, state transitions, and safety checks. This separation not only improves maintainability but also ensures that logic, configuration, and oversight remain clearly distinct.

The schema is further optimized for flexibility and performance through the use of JSONB fields, GIN indexes, and array structures, allowing complex rules and state data to be represented in a compact and query-efficient manner. Real-time responsiveness is achieved with database triggers and the pg\_notify mechanism, which broadcast updates to critical control elements as they occur. At the same time, embedded audit logging ensures that every operational decision and safety-relevant change is captured with full traceability, meeting the demands of regulatory oversight. By combining circuit-based occupancy tracking with these advanced features, the schema provides both the precision needed for safe operations and the extensibility required for long-term evolution.